\clearpage
\SourceLeaf{189}
\thispagestyle{empty}
\pdfbookmark[0]{Reviews of Other Authors' Writings}{other-notices-title}
\begin{center}
\vspace*{1in}{\Large REVIEWS\par}
\vspace{.4in}{\large OF OTHER AUTHORS'\par}
\vspace{.4in}{\Large WRITINGS.\par}
\end{center}

\clearpage
\SourceLeaf{191}
\renewcommand{\BookRunningTitle}{Reviews of Other Authors' Writings}
\pdfbookmark[1]{Recherches sur l'irréductibilité. Review, 1809, March 11}{notice-1809-03-11}
\begin{center}
\rule{.86\linewidth}{.3pt}\par
Göttingische gelehrte Anzeigen.\quad 1809, March 11.\par
\rule{.86\linewidth}{.3pt}
\end{center}
\begin{center}
\emph{Recherches sur l'irréductibilité Arithmétique et Géométrique des nombres et de leurs puissances.}\quad 1808.\quad (No place of publication.\quad 25 pages in large quarto.)
\end{center}

A work whose aim is to represent irrational radical quantities in the form of rational quantities. We must content ourselves with having drawn the attention of readers interested in mathematics to this little work, since the limits of these pages do not allow us to give a fuller account and examination here of the author's distinctive standpoint and treatment of radical quantities, which departs entirely from the usual one.

\Needspace{8\baselineskip}
\pdfbookmark[1]{Chernac, Cribrum Arithmeticum. Review, 1812, March 23}{notice-1812-03-23}
\begin{center}
\rule{.86\linewidth}{.3pt}\par
Göttingische Gelehrte Anzeigen.\quad 1812, March 23.\par
\rule{.86\linewidth}{.3pt}
\end{center}
\emph{Cribrum Arithmeticum, sive tabula continens numeros primos a compositis segregatos, occurrentes in serie numerorum ab unitate progredientium usque ad decies centena milia et ultra haec ad viginti millia} $(1020000)$. \emph{Numeris compositis, per} $2$, $3$, $5$ \emph{non dividuis, adscripti sunt divisores simplices, non minimi tantum, sed omnino}
\SourceLeaf{192}
\emph{omnes. Confecit} \textsc{Ladislaus Chernac}, \emph{Pannonius, A. L. M. Philos. et Medic. Doctor, in almo lyceo Daventriensi philosophiae professor. Daventriae} 1811. (At the author's expense, printed by J. H. Lange.\quad XXII and 1022 pages, large quarto.)

The full title of this important and highly meritorious work already sufficiently indicates its contents: it is a table of all the prime factors of all numbers not divisible by $2$, $3$ and $5$, from $1$ to $1020000$, calculated through several years of work as careful as it was laborious, and printed neatly and, as far as we have found by occasional checks, very correctly. Anyone who frequently carries out calculations with larger numbers can readily judge the value of such a gift to arithmetic. The author deserves double thanks: both for the exceedingly laborious work itself, by which he has placed his name beside the unforgettable names of \textsc{Rhaeticus}, \textsc{Pitiscus}, \textsc{Brigg}, \textsc{Vlacq}, \textsc{Wolfram}, \textsc{Taylor} and others, and for the certainly very considerable expense incurred in printing it, for which a publisher would otherwise scarcely have been found. Such tables have often been calculated before, though mostly on a smaller scale, but have either remained entirely in manuscript or have not been completely printed. As is well known, \textsc{Lambert} formerly did all he could to encourage the continuation of \textsc{Pell}'s often reprinted table, which extends to $100000$. According to a report by \textsc{Bernoulli} in \textsc{Lambert}'s correspondence, \textsc{Oberreit} had extended it to $500000$, and a copy had come into \textsc{Schulze}'s hands. \textsc{Anton Felkel}, as reported in the Monatl. Correspondenz, volume 2, page 223, had completed it in manuscript up to two million and intended later to provide it up to $2460000$; but what had already been printed in Vienna at public expense was used for cartridges in the Turkish war because no buyers could be found! Thus many years of meritorious work were lost to the public: all the more, therefore, we considered it our duty to announce the publication of the present work here. The first million is now available for everyone's use; anyone with the opportunity and enthusiasm for this subject may therefore direct their efforts to what follows.

\SourceLeaf{193}
\Needspace{8\baselineskip}
\pdfbookmark[1]{Burckhardt, deuxième million. Review, 1814, November 3}{notice-1814-11-03}
\begin{center}
\rule{.86\linewidth}{.3pt}\par
Göttingische gelehrte Anzeigen.\quad 1814, November 3.\par
\rule{.86\linewidth}{.3pt}
\end{center}
\emph{Tables des diviseurs pour tous les nombres du deuxième million, ou plus exactement depuis} $1020000$ \emph{à} $2028000$, \emph{avec les nombres premiers qui s'y trouvent. Par} \textsc{J. Ch. Burckhardt}, \emph{membre de l'institut impérial, du bureau des longitudes de France, et de plusieurs autres sociétés savantes. Paris}, 1814. \emph{M\textsuperscript{me} V\textsuperscript{e} Courcier}. (VIII. and 112 pages in folio.)

Sooner than we would have dared to hope when reviewing \textsc{Chernac}'s factor table covering the first million, we can already report the completion and publication of a similar table for the second million. Through this laborious work the distinguished author, whose name itself guarantees the greatest care and accuracy, has placed all those interested in arithmetic deeply in his debt. \textsc{Chernac}'s table for the first million gives all the prime factors; \textsc{Burckhardt}'s for the second, however, gives only the least divisor in each case. The complete factorization of a number in the second million therefore requires division by the least divisor and looking up the quotient in \textsc{Chernac}'s table. But this small effort is of no significance compared with the great advantage of having the table occupy so much less space, while the prospect remains of seeing it extended in time to ten million. The author made it possible to fit the table into this small volume partly by restricting it to the least divisor and partly by printing it as economically as possible. If $a$ stands for any of the eighty numbers below $300$ that are not divisible by $2$, $3$ and $5$, then every number not divisible by $2$, $3$ and $5$ is included in the form $300n+a$. All eighty numbers for which $n$ has the same value are found in a single vertical column, and each page contains thirty such columns. Thus each page includes, out of nine thousand numbers proceeding in their natural order, all those not divisible by $2$, $3$ or $5$.

\Needspace{6\baselineskip}
The method by which Mr. \textsc{Burckhardt} constructed his table deserves special mention here. He had a grid engraved on copper, in which $81$ horizontal and $78$ vertical lines formed a rectangle divided into $80\times77$, that is, $6160$ small squares, and had the necessary number of impressions made from it. At the side, the eighty values
\SourceLeaf{194}
of $a$ could also be engraved immediately; the values of $300n$, in continuous order, were written with a pen above the $77$ vertical columns. Thus each sheet represents all the numbers not divisible by $2$, $3$ and $5$ found among each set of $23100$ numbers proceeding in their natural order, and $44$ sheets suffice to cover a whole million. It is easy to see that the numbers whose least divisor is $7$ or $11$ recur in the same order on every following sheet. These divisors could therefore be engraved directly on the copperplate and consequently appeared automatically in their proper places on every sheet. To enter the following divisors, for example $13$, Mr. B. took only $13$ columns across the width of a spare sheet and, treating it as the beginning of his table, cut out all the squares that had to contain the divisor $13$. He therefore needed only to lay this stencil over the first thirteen columns of the first sheet, then over the next thirteen, and so on, to see at once all the places that had to be filled, insofar as they did not already contain $7$ or $11$, with $13$. The same procedure was then used for the divisor $17$, and so on. The spare sheets sufficed in this way up to the divisor $73$; for the larger divisors $79$, $83$, and so on, Mr. B. appears to have assembled the frame from two or more pieces. For divisors exceeding $500$, however, Mr. B. preferred to find the multiples by addition, needing to take only primes for the other factor. We find this whole procedure highly suitable, and would recommend it for imitation to anyone inclined to extend the table further. For the third and fourth millions, the author himself has meanwhile already carried out a large part of the calculations, so we have good reason to hope to see these too made available in print soon.

\Needspace{22\baselineskip}
\pdfbookmark[1]{Burckhardt, troisième million. Review, 1816, November 7}{notice-1816-11-07}
\begin{center}
\rule{.86\linewidth}{.3pt}\par
Göttingische gelehrte Anzeigen.\quad 1816, November 7.\par
\rule{.86\linewidth}{.3pt}
\end{center}
\emph{Tables des diviseurs pour tous les nombres du troisième million, ou plus exactement, depuis} $2028000$ \emph{à} $3036000$, \emph{avec les nombres premiers qui s'y trouvent; par} \textsc{J. Chr. Burckhardt}, \emph{membre de l'académie royale des sciences, du bureau des }
\SourceLeaf{195}
\emph{longitudes de France et de plusieurs autres sociétés savantes. Paris} 1816. \emph{M\textsuperscript{me} V\textsuperscript{e} Courcier}. (112 pages in folio.)

Since in our review of the table of factors for the second million we have already described in detail both the method of calculation used by the distinguished author and the arrangement of the table itself, we can here content ourselves with announcing the publication of the table for the third million. Shortly we may also expect from the author the table for the \emph{first} million, presented in the same way, so that the whole table extending beyond three million will then occupy only a moderately sized volume. For this the author deserves the thanks of all those interested in arithmetic, who see this laborious work meeting a need on a scale far exceeding anything one could have dared to hope for only a few years ago.

\Needspace{8\baselineskip}
\pdfbookmark[1]{Burckhardt, premier million. Review, 1817, August 9}{notice-1817-08-09}
\begin{center}
\rule{.86\linewidth}{.3pt}\par
Göttingische gelehrte Anzeigen.\quad 1817, August 9.\par
\rule{.86\linewidth}{.3pt}
\end{center}
\emph{Tables des diviseurs, pour tous les nombres du premier million, ou plus exactement depuis} $1$ \emph{à} $1020000$, \emph{avec les nombres premiers qui s'y trouvent; par} \textsc{J. Chr. Burckhardt}, \emph{membre de l'académie des sciences dans l'institut royal, du bureau des longitudes de France, et de plusieurs autres sociétés savantes. Paris} 1817. \emph{M\textsuperscript{me} V\textsuperscript{e} Courcier}. (114 pages in folio.)

Referring here to our reviews of the tables for the second and third millions, we now merely announce the actual\textperiodcentered\ publication of these factor tables for the first million. We thus now possess a connected whole for the first three millions. For the present first million, the author used partly \textsc{Chernac}'s \emph{Cribrum Arithmeticum} and partly a manuscript table by \textsc{Schenmark} held by the library of the Royal Institute. The latter had not, however, been constructed with quite all the care one could wish for, and deciding which was right in cases where the two disagreed was often rather laborious. Only a very small number of errors were found in \textsc{Chernac}'s table, which Mr. \textsc{Burckhardt} has reported here.

\SourceLeaf{196}
The author also already has most of the materials ready for the fourth, fifth and sixth millions, and offers to supply this continuation if the publisher is encouraged by sufficient sales of the first three millions. It would indeed be very regrettable if the fruits of such laborious and useful work were withheld from the world.
